\documentclass[11pt]{article}
\usepackage[T1]{fontenc}
\usepackage[utf8]{inputenc}
\usepackage{lmodern}
\usepackage[margin=1in]{geometry}
\usepackage{array}
\usepackage{booktabs}
\usepackage{enumitem}
\usepackage{longtable}
\usepackage{tabularx}
\PassOptionsToPackage{hyphens}{url}
\usepackage[hidelinks]{hyperref}
\usepackage{xurl}
\setlength{\parindent}{0pt}
\setlength{\parskip}{0.6em}
\setlist[itemize]{leftmargin=1.4em}
\setlist[enumerate]{leftmargin=1.6em}
\emergencystretch=2em
\hbadness=10000
\sloppy

\title{ISO/IEC/IEEE 42010:2022\\Citation-Clean Handout}
\author{}
\date{}

\begin{document}
\maketitle

\section*{Purpose}
This handout provides a concise, source-grounded summary of ISO/IEC/IEEE 42010:2022 for practitioners who need a clean explanation of what the standard covers, what it requires, and what it deliberately leaves open.

\section*{What the Standard Is}
ISO/IEC/IEEE 42010:2022 is the current international standard for \emph{architecture description}. The official ISO record states that it specifies requirements for the structure and expression of an architecture description (AD) for a wide range of entities, including software, systems, enterprises, systems of systems, families of systems, products, product lines, service lines, technologies, and business domains \cite{iso42010}.

The standard distinguishes the \emph{architecture} of an entity of interest from the \emph{architecture description} that expresses that architecture. In the ISO abstract, architectures themselves are explicitly not the subject of the document; the document governs how an AD is structured and expressed \cite{iso42010}.

\section*{Core Concepts}
\begin{itemize}
    \item \textbf{Entity of interest.} The subject whose architecture is being described. The 2022 edition uses the term \emph{entity of interest} rather than the older \emph{system of interest} terminology \cite{iso42010,nora42010}.
    \item \textbf{Stakeholders and concerns.} Architecture descriptions are created to address the concerns of stakeholders. The 42010 conceptual model places stakeholders and concerns at the center of why views are created \cite{cm42010}.
    \item \textbf{Viewpoints and views.} A viewpoint specifies conventions for constructing a view, and a view is the resulting representation of the architecture from that viewpoint \cite{cm42010}.
    \item \textbf{Model kinds.} The ISO abstract states that the standard specifies requirements for model kinds as part of supporting useful architecture descriptions \cite{iso42010}.
    \item \textbf{Architecture description frameworks and languages.} The ISO abstract also states that the standard specifies requirements for an architecture description framework (ADF) and an architecture description language (ADL) \cite{iso42010}.
    \item \textbf{Correspondence and consistency.} The conceptual model includes correspondences among elements in an architecture description so that relationships and consistency can be expressed across views and models \cite{cm42010}.
    \item \textbf{Architecture decisions and rationale.} The standard framework emphasizes that an AD should not only show structure but also preserve the rationale behind significant architectural choices so that the description is understandable, reviewable, and maintainable over time \cite{cm42010}.
\end{itemize}

\section*{What the 2022 Edition Clarifies or Updates}
The strongest edition-level changes that are easy to support from current public sources are these:

\begin{itemize}
    \item The standard is now explicitly titled \emph{Software, systems and enterprise --- Architecture description} \cite{iso42010}.
    \item The terminology shifts from \emph{system of interest} to \emph{entity of interest} \cite{nora42010}.
    \item \emph{Architecture description framework} is used in place of the older \emph{architecture framework} terminology \cite{nora42010}.
    \item A viewpoint may govern one or more architecture views in an architecture description \cite{nora42010}.
    \item Model kinds are treated as an explicit conformance-related concept in the 2022 edition \cite{iso42010,nora42010}.
\end{itemize}

\section*{What the Standard Does Not Prescribe}
The ISO abstract is clear that ISO/IEC/IEEE 42010:2022 does \emph{not} specify:
\begin{itemize}
    \item the architecting process to use,
    \item the methods used to develop the architecture,
    \item the modeling notations, techniques, or tools to use,
    \item the medium or recording format for the AD,
    \item or the actual characteristics of the entity being described \cite{iso42010}.
\end{itemize}

\section*{Practical Reading of the Standard}
For working architects, the standard is best read as a \textbf{content and traceability discipline}. It does not tell you to use ArchiMate, UML, BPMN, SysML, spreadsheets, prose, or a repository tool. Instead, it expects you to make those choices explicit, tie them to stakeholder concerns, and organize the resulting architecture description so that a reviewer can understand what concerns were addressed, by which viewpoints, in which views, with what rationale and correspondences \cite{iso42010,cm42010}.

\section*{What a Conforming Architecture Description Usually Needs to Show}
A practical architecture description aligned to ISO 42010 typically needs to make the following visible:

\begin{center}
\renewcommand{\arraystretch}{1.2}
\begin{tabularx}{\textwidth}{>{\raggedright\arraybackslash}p{0.28\textwidth} X}
\toprule
\textbf{Element} & \textbf{Why it matters} \\
\midrule
Entity of interest and scope & Establishes what is being described and what is outside scope. \\
Stakeholders and concerns & Explains why the architecture description exists and what it must answer. \\
Selected viewpoints & Shows the rules used to construct views. \\
Architecture views & Provides the actual architectural representations used for communication and review. \\
Model kinds and languages & Makes modeling conventions explicit. \\
Decisions and rationale & Preserves traceability for major architectural choices. \\
Correspondences & Shows how elements across views relate and stay consistent. \\
Framework or organizing scheme & Provides structure for how the AD is assembled and maintained. \\
\bottomrule
\end{tabularx}
\end{center}

\section*{Bottom Line}
ISO/IEC/IEEE 42010:2022 is not a method, notation, or tool standard. It is a standard for creating architecture descriptions that are explicit about \emph{who} they serve, \emph{which concerns} they address, \emph{which viewpoints and views} they use, and \emph{how the pieces relate}. That is why it remains relevant across software architecture, systems engineering, and enterprise architecture \cite{iso42010,cm42010}.

\section*{References}
\begin{thebibliography}{9}
\bibitem{iso42010}
International Organization for Standardization (ISO), \emph{ISO/IEC/IEEE 42010:2022 - Software, systems and enterprise --- Architecture description}, official standard record and abstract. Available at: \url{https://www.iso.org/standard/74393.html}

\bibitem{cm42010}
ISO/IEC/IEEE 42010 community site, \emph{ISO/IEC/IEEE 42010: Conceptual Model}. Available at: \url{https://www.iso-architecture.org/42010/cm/}

\bibitem{nora42010}
NORA Online, \emph{ISO 42010}. Used here only as a public summary of terminology and framing differences relevant to the 2022 edition. Available at: \url{https://www.noraonline.nl/wiki/ISO_42010}
\end{thebibliography}

\end{document}



